\documentclass[11pt]{article}
\usepackage[margin=1in]{geometry}
\usepackage{amsmath,amssymb}
\usepackage{booktabs}
\usepackage{array}
\usepackage{graphicx}
\usepackage[hidelinks]{hyperref}
\usepackage{enumitem}

\title{\textbf{Comparing Exploration Algorithms in the 3D Maze}\\
\large A Course-Project Reproduction under the Paper's Appendix~C.2 Settings\\
\normalsize \emph{Challenging Problems in Reinforcement Learning} (SS~2026)\\
\normalsize Topic: Exploration vs.\ Exploitation}
\author{}
\date{2026-05-31\\ \normalsize\emph{results updated 2026-06-03: LPM corrected to the paper's
paper-faithful \emph{log-space} reward (Eq.~3); Appendix-C.2 config ($\lambda_i{=}1$, entropy
$0.03$, raw reward); 6 policies $\times$ 3 variants $\times$ 64 seeds}}

\newcommand{\cov}{C}

\begin{document}
\maketitle

\begin{abstract}
\noindent We compare five exploration mechanisms---Learning Progress Monitoring (LPM), Random
Network Distillation (RND), the Intrinsic Curiosity Module (ICM), a raw next-state prediction-error
bonus (MSE), and a no-intrinsic A2C baseline (\textsc{none})---plus a \emph{uniform-random control}
(no learning at all) in the Miniworld 3D maze of Hou, An \& Du (2026), \emph{Beyond Noisy-TVs:
Noise-Robust Exploration via Learning Progress Monitoring}. The maze emits \emph{no extrinsic
reward}, so a method's quality is its \emph{spatial coverage}, and the noisy-TV failure mode is
geometric: a random-RGB ``noise wall'' at $z\approx 8$ separates the two lower spawn rooms from the
far room. \textbf{This version corrects LPM to the paper's faithful \emph{log-space} reward}
$r^i=g_\phi-\log\mathrm{MSE}$ (Eq.~3), replacing the released notebooks' raw-error clipped form; doing
so also required lowering the error-model learning rate (the notebook value silently broke the
log-space objective). We run the paper's exact Appendix~C.2 configuration ($\lambda_i{=}1$, entropy
$0.03$, raw intrinsic reward) at the paper's budget but with $64$ seeds ($5\times10^{4}$ steps each).
Three findings stand out. \textbf{(1)} The paper's central claim---\emph{noise robustness}---reproduces
cleanly with the faithful reward: under action-triggered noise, LPM selects the noisy-TV action only
$3.0\%$ of the time (median $0.9\%$; $1/64$ seeds) while the raw-pixel method (MSE) is trapped on it
$83\%$ of the time and its coverage collapses. \textbf{(2)} On \emph{coverage}, no learned method
beats the pure uniform-random control: uniform covers the most in \emph{all three} variants. In this
small maze coverage rewards undirected randomness, so the paper's ``LPM achieves the highest
coverage'' ranking does \emph{not} reproduce, while its noise-robustness claim---decoupled from
coverage---does. \textbf{(3)} The faithful, \emph{unclipped} $\lambda_i{=}1$ reward is noise-robust but
\emph{unstable}: $45\%$ of LPM's \texttt{action\_noise} seeds collapse to degenerate near-stationary
policies, pushing LPM's coverage below ICM in both mean and median---a reversal from the earlier
clipped run. The collapse is policy instability, \emph{not} TV-fixation (the collapsed seeds still
avoid the TV), and the notebook's clip was masking it. A $64$-seed-averaged, log-scaled
occupancy-density figure renders the spatial dynamics.
\end{abstract}

\section{Motivation and research questions}
The course topic is the exploration--exploitation trade-off. Deep RL drives directed exploration
with intrinsic-motivation bonuses (prediction error, novelty, learning progress). The \emph{noisy-TV}
problem is the canonical failure mode of prediction-error exploration: an unlearnable source of
randomness produces perpetually high prediction error, so a curiosity-driven agent fixates on it.
LPM rewards the \emph{improvement} of its dynamics model rather than the error level, and is
therefore noise-robust by construction. The 3D maze makes this failure spatial and directly
visualisable: the noise wall is a physical place the agent can get stuck at.

The maze geometry is a hand-designed four-room layout in an $18\times 12$ world (continuous units):
two spawn rooms below the wall (room1 $x\!\in\![0,4]$, room2 $x\!\in\![14,18]$, both $z\!\in\![0,8]$),
the thin grass-textured \emph{noise wall} (room3, $z\!\in\![8,8.1]$), and the large far room
(room4, $z\!\in\![8.1,12]$). The agent spawns in room1 and must traverse the wall to explore room4.

\begin{description}[leftmargin=2.2em,style=nextline]
  \item[RQ1 (coverage \& noise robustness, reproduction).] Under the noisy variants
  (\texttt{noisy\_tv}: random-RGB wall $+$ 25\% sticky actions; \texttt{action\_noise}: an action
  that shows a random CIFAR-10 image) vs.\ the clean control (\texttt{nonoise}), which methods keep
  expanding coverage and which stall---and crucially, do any of them beat a \emph{uniform-random}
  policy that does no learning at all?
  \item[RQ2 (spatial dynamics, novelty).] How does the spatial distribution of exploration evolve
  over training? We answer this with a coverage-heatmap-evolution figure (rows: method; columns:
  training-time window) \emph{averaged over all 64 seeds}.
\end{description}

\section{Methods}
All learned methods use the same A2C base, configured to the paper's Appendix~C.2 (MiniWorld):
RMSprop, policy LR $0.01$, $\gamma=0.99$, GAE $\lambda=0.95$, value-loss coef.\ $0.5$, gradient
clipping $0.5$, a single-epoch update every $64$ steps, dynamics-model LR $10^{-3}$, error-buffer
size $100$, entropy coefficient $0.03$, and intrinsic weight $\lambda_i=1$, over a Nature-style CNN
on $160\times120$ RGB frames. Because the maze extrinsic reward is identically zero, the per-step
training reward is the \emph{raw} intrinsic term,
\begin{equation}
  r_t \;=\; r^{\text{ext}}_t + \lambda_i\, r^{\,i}_t \;=\; \lambda_i\, r^{\,i}_t,
  \qquad \lambda_i = 1,
  \label{eq:reward}
\end{equation}
with no running-std normalisation. (The upstream notebooks instead used $\lambda_i=0.1$, entropy
$0.05$, and a normalised reward---settings that \emph{disagree} with the paper's own Appendix~C.2
and make the intrinsic signal $10\times$ too weak; we follow C.2. The deviation is logged in
\texttt{UPSTREAM.md}.) For \textsc{none}, $r_t\equiv 0$, so the A2C policy is still \emph{trained}
(by the entropy bonus and value bootstrap) but receives no reward signal; the \emph{uniform} control
is not trained at all---it samples actions $a\sim\mathrm{Uniform}\{0,\dots,4\}$ every step.

\paragraph{LPM corrected to the paper's log-space reward.} The released maze notebooks (and our
earlier runs) computed the LPM bonus in \emph{raw}-MSE space with a clip,
$r^i=\min\!\big(0.5,\ \eta\,e^{g_\phi}-\mathrm{MSE}\big)$ with $\eta{=}1$, which does \emph{not} match
the paper's Eq.~(1)--(3). We corrected LPM to the faithful log-space form: with the per-step error
defined as $\varepsilon=\log\mathrm{MSE}$ (Eq.~1) and an error model $g_\phi$ trained to predict the
\emph{previous} dynamics model's log-error $\mathbb{E}_{\mathcal D}[\varepsilon^{(\tau-1)}]$ (Eq.~2),
the intrinsic reward is their difference (Eq.~3),
\begin{equation}
  r^{\,i}_t \;=\; g_\phi(o_t,a_t)\;-\;\varepsilon^{(\tau)}_t
  \;=\; \mathbb{E}_{\mathcal D}\!\big[\log\mathrm{MSE}_{\text{prev}}\big]\;-\;\log\mathrm{MSE}_{\text{cur}},
  \label{eq:lpm}
\end{equation}
with no $\eta$ and no clip, gated to $0$ until the error buffer fills ($|\mathcal D|{=}100$, Alg.~1),
and $g_\phi$ updated together with the dynamics model every step. This also required lowering the
error-model learning rate to $10^{-3}$: under the log-space objective the notebook's $10^{-2}$ drives
$g_\phi$ into its $[-10,10]$ clamp (zero gradient---a dead error model that pins the reward at a
constant $\approx-6$); at $10^{-3}$ the error model fits the log-error targets and the bonus becomes
the intended learning-progress signal, fluctuating around $0$ and decaying as the model saturates
(cf.\ the paper's Fig.~2). The pre-fix raw form is preserved behind \texttt{reward\_space="raw"}; the
fix and its root-cause are logged in \texttt{UPSTREAM.md}.

\begin{table}[h]
\centering
\begin{tabular}{@{}llp{6.4cm}@{}}
\toprule
\textbf{Class} & \textbf{Method} & \textbf{Mechanism (intrinsic reward $r^{i}$)} \\
\midrule
Learning progress        & \textbf{LPM} & $g_\phi-\varepsilon=\mathbb{E}_{\mathcal D}[\log\mathrm{MSE}_{\text{prev}}]-\log\mathrm{MSE}_{\text{cur}}$ (Eq.~\ref{eq:lpm}): expected-minus-actual \emph{log}-error (model improvement) \\
Prediction error         & RND          & error of a predictor vs.\ a fixed random target network on the next frame \\
Prediction error (feat.) & ICM          & forward-dynamics error in an inverse-dynamics feature space \\
Prediction error (pixel) & MSE          & raw next-state pixel-prediction error (the pure noisy-TV victim) \\
None (baseline)          & \textsc{none}& no intrinsic reward; A2C \emph{is} still trained (entropy + value bootstrap) \\
Control                  & uniform      & $a\sim\mathrm{Uniform}\{0\dots4\}$; \emph{no learning} (pure random walk) \\
\bottomrule
\end{tabular}
\caption{The six policies. LPM and MSE share the same pixel-decoder predictor, isolating the effect
of the learning-progress signal vs.\ raw error. ICM and RND are canonical implementations added on
the shared CNN encoder. The \emph{uniform} control is the proper ``does any of this beat chance?''
baseline---distinct from \textsc{none}, whose policy is still shaped by training.}
\end{table}

\section{Experimental design}
A single manipulation grid: policy $\times$ variant $\times$ seed.

\paragraph{Grid.} Six policies (LPM, RND, ICM, MSE, \textsc{none}, uniform) $\times$ three variants
(\texttt{nonoise}, \texttt{noisy\_tv}, \texttt{action\_noise}) $\times$ \textbf{64 seeds} $=1152$
runs. Each run is a single environment trained for $5\times10^{4}$ steps---the paper's maze budget
(Fig.~3 of Hou et al.\ trains for $5\times10^{4}$ exploration steps). We use $64$ seeds (the paper
used $10$) to tighten the estimates and resolve LPM's bimodal seed behaviour.

\paragraph{Budget.} The runs are independent, so the grid is embarrassingly parallel: on a 128-core
Linux box, \texttt{run\_grid.py --jobs 120 --threads-per-job 1} finished the full $1152$-run grid in
$\approx\!6$\,h at $\sim\!$full CPU ($0$ failures; peak $\sim\!140$\,GB of 1.1\,TB RAM, set by the
$120$-way concurrency). Single-environment A2C is a batch-size-1 forward pass, so one CPU thread per
job is optimal (two threads gave only $1.25\times$ per run)---the speed-up comes from running many
processes. Runs use CPU (the box has no GPU); the $160\times120$ frames render headless via Mesa EGL
(\texttt{PYGLET\_HEADLESS=true}, \texttt{LP\_NUM\_THREADS=1}). The grid is resumable---a run's
position \texttt{.npz} marks it complete.

\section{Metrics}
Coverage is measured on the maze's native $72\times48$ occupancy grid (cells of $0.25$ world units),
restricted to the \emph{reachable} cells (the union of the four rooms; the empty corridor between the
spawn rooms is excluded). Let the final value of a quantity be its mean over the last $10\%$ of
logged updates, reported as mean$\,\pm\,$std over the 64 seeds.

\begin{itemize}[leftmargin=1.4em]
  \item \textbf{Total coverage:} $\cov = (\text{visited reachable cells})/(\text{reachable cells})$,
  reported both as a final number and as a learning curve vs.\ frames.
  \item \textbf{TV-action share:} fraction of steps spent on the noisy-TV action (action~4) under
  \texttt{action\_noise}---the direct measure of noisy-TV fixation. Prediction: highest for the pixel
  prediction-error method (MSE), lowest for LPM, exactly $1/5$ for the uniform control.
  \item \textbf{Coverage-heatmap evolution (averaged):} for each variant, a
  rows(method)$\times$columns(training-window) grid of top-down occupancy heatmaps \emph{averaged
  over all 64 seeds}. The primary view is \emph{occupancy density} = the mean per-cell visit count,
  drawn on a single \emph{shared log colour scale} across all panels (a colorbar is overlaid). The
  raw per-cell count spans $\sim\!0.1$ to $\sim\!10^{3}$---a few dwell/spin cells are visited
  $1000\times$ while typical cells see $1$--$5$ visits---so a per-panel \emph{linear} scale (our
  earlier figures) washed the path structure out; the shared log scale exposes both \emph{where} each
  policy dwells and how much. A companion \emph{coverage-frontier} view (fraction of seeds that ever
  visited each cell, $0$--$1$) is also produced. The maze rooms and the $z=8$ wall line are overlaid.
\end{itemize}

\section{Results}
We ran the full $1152$-run grid at $5\times10^{4}$ steps under the C.2 configuration with the
corrected log-space LPM. The single-environment A2C is high-variance (between-seed std
$0.08$--$0.32$), so we flag which conclusions are robust.

\paragraph{Noise robustness reproduces---cleanly, even with the faithful reward.} In
\texttt{action\_noise}, action~4 replaces the observation with a random CIFAR-10 image. The corrected
log-space LPM selects it on only $3.0\%\pm11.3\%$ of steps (\emph{median} $0.9\%$; only $1$ of $64$
seeds ever fixates), because its dynamics model cannot \emph{improve} on the unlearnable image, while
the raw-pixel predictor MSE is trapped on it $83\%\pm27\%$ of the time and its coverage collapses to
$0.13$ (Table~\ref{tab:tv}, Fig.~\ref{fig:tv}). The curiosity baselines are partially fooled
(RND $67\%$, ICM $53\%$). The uniform control sits at exactly $0.200\pm0.003$ (the $1/5$ sanity
check). This is the paper's central claim and it reproduces.

\begin{table}[h]
\centering
\begin{tabular}{@{}lc@{}}
\toprule
\textbf{Policy} & \textbf{TV-action share} (\texttt{action\_noise}) \\
\midrule
\textbf{LPM}  & $\mathbf{0.030 \pm 0.113}$ \\
uniform       & $0.200 \pm 0.003$ \\
\textsc{none} & $0.205 \pm 0.175$ \\
ICM           & $0.526 \pm 0.242$ \\
RND           & $0.666 \pm 0.276$ \\
\textbf{MSE}  & $\mathbf{0.830 \pm 0.266}$ \\
\bottomrule
\end{tabular}
\caption{Noisy-TV fixation under \texttt{action\_noise}, mean$\,\pm\,$std over 64 seeds. LPM
essentially never stares at the noisy TV ($3.0\%$ mean, $0.9\%$ median, $1/64$ seeds fixate); MSE is
trapped on it ($83\%$); the curiosity baselines are partially fooled; the uniform control is at
chance ($1/5$).}
\label{tab:tv}
\end{table}

\paragraph{Coverage: no learned method beats a uniform-random policy.} The headline coverage result
is negative and important (Table~\ref{tab:cov}, Fig.~\ref{fig:curves}): the \emph{uniform-random}
control covers the most in \emph{all three} variants---and with by far the lowest variance (std
$0.08$--$0.13$ vs $0.19$--$0.32$ for the learned methods); in clean \texttt{nonoise} uniform ($0.587$)
now edges even the best learned method (ICM $0.542$). No learned policy, LPM included, exceeds the
no-learning control anywhere. Training \emph{reduces} coverage here: any learned policy
(intrinsic-driven or the trained \textsc{none}) commits to a narrower, repetitive walk, whereas
uniform thrashing maximally spreads through a small maze. Coverage in this setting therefore rewards
\emph{randomness}, not directed exploration, and the paper's ``LPM achieves the highest coverage''
ranking does \emph{not} reproduce. Note that noise-avoidance and coverage are \emph{decoupled}: under
\texttt{action\_noise} the uniform policy wastes $20\%$ of its steps frozen on the TV yet still covers
most ($0.696$), while LPM almost never touches the TV ($3.0\%$) yet covers far less on average
($0.299$). (Uniform covers \emph{more} under the noisy variants than under clean \texttt{nonoise}
because those variants add $25\%$ sticky actions, which create forward momentum and carry a random
walker further.)

\paragraph{The faithful $\lambda_i{=}1$ reward is noise-robust but \emph{unstable}.} With the clip
removed, the log-space bonus at $\lambda_i{=}1$ is strong, and LPM's behaviour is sharply
\emph{bimodal}: under \texttt{action\_noise}, $29$ of $64$ seeds ($45\%$) collapse to a degenerate
near-stationary policy ($<\!0.05$ coverage) while the rest reach $0.3$--$0.8$. This pulls LPM's
coverage \emph{below} ICM in both mean ($0.299$ vs $0.410$) \emph{and} median ($0.101$ vs $0.416$)---a
reversal from the earlier clipped-reward run, where the clip damped the reward and fewer seeds
collapsed. Crucially this is \emph{policy instability, not noisy-TV susceptibility}: the collapsed LPM
seeds freeze in a tight loop \emph{somewhere in the maze}, not on the TV (LPM's TV-share stays at
$3\%$, versus ICM's $53\%$), so the collapse lowers coverage without breaking noise-robustness. The
clip in the released notebooks was therefore silently \emph{masking} a stability problem of the
unclipped $\lambda_i{=}1$ reward; the faithful formula exposes it. This is a concrete
stability-ablation hook (reward scaling $\lambda_i$, reward/return normalisation, or an advantage
clip), and it reinforces that coverage---not the metric LPM is built for---is the wrong yardstick in
this maze.

\begin{table}[h]
\centering
\begin{tabular}{@{}lccc@{}}
\toprule
\textbf{Policy} & \texttt{nonoise} & \texttt{noisy\_tv} & \texttt{action\_noise} \\
\midrule
\textbf{uniform} & $\mathbf{0.587\,(0.13)}$ & $\mathbf{0.710\,(0.08)}$ & $\mathbf{0.696\,(0.09)}$ \\
\textsc{none}    & $0.459\,(0.24)$ & $0.611\,(0.25)$ & $0.589\,(0.25)$ \\
LPM              & $0.332\,(0.25)$ & $0.452\,(0.27)$ & $0.299\,(0.32)$ \\
MSE              & $0.476\,(0.23)$ & $0.551\,(0.26)$ & $0.129\,(0.12)$ \\
ICM              & $0.542\,(0.20)$ & $0.536\,(0.29)$ & $0.410\,(0.19)$ \\
RND              & $0.468\,(0.22)$ & $0.519\,(0.29)$ & $0.290\,(0.22)$ \\
\bottomrule
\end{tabular}
\caption{Coverage fraction, mean (std) over 64 seeds, under the C.2 configuration with the corrected
log-space LPM. Best per column in bold: the uniform-random control wins \emph{all three} variants.
Every learned method falls below the no-learning control. LPM's coverage is depressed by a bimodal
collapse (29/64 seeds under \texttt{action\_noise}; median $0.101$); the collapse is policy
instability, not TV-fixation (see text and Table~\ref{tab:tv}).}
\label{tab:cov}
\end{table}

\paragraph{Why the coverage ranking is not reproducible here (structural).}
\begin{enumerate}[leftmargin=1.6em]
  \item In a small maze ($\sim\!1500$ reachable cells) with no extrinsic reward, a random walk
  already covers most cells, so there is nothing for directed exploration to add; committing to any
  policy only narrows the walk. Intrinsic motivation targets \emph{sparse-reward, hard-exploration}
  tasks (e.g.\ Montezuma's Revenge), not pure coverage of a small room.
  \item The maze notebooks contain only LPM and pixel-curiosity---\emph{none} of Figure~3's baselines
  (AMA, EDT, EME, Ensemble, IDF, RND). Figure~3 was produced by a fuller pipeline not in the released
  maze code; our RND/ICM are our own additions on the shared encoder.
  \item Single-environment A2C at $5\times10^{4}$ steps is dominated by between-seed variance
  (std $\approx$ mean) and the curves have not plateaued at $50$k; for LPM the variance is genuinely
  bimodal (collapse vs.\ explore) rather than merely noisy.
\end{enumerate}

\paragraph{Spatial dynamics (RQ2).} The 64-seed-averaged occupancy-density figure
(Fig.~\ref{fig:heat}, \texttt{action\_noise}) makes the story legible once the shared log colour
scale is used: brightness is the mean per-cell visit count. The \emph{uniform} row spreads broadly
and fairly evenly across both spawn rooms and up past the wall, at moderate per-cell counts; the
trained policies instead \emph{concentrate}---\textsc{none} and especially MSE show bright
($10^{2}$--$10^{3}$) dwell/spin cells where the agent loops in place, with sparser coverage
elsewhere, while LPM/ICM/RND are patchier and dimmer (the collapsed LPM seeds contribute tight,
bright single-cell loops that average into faint hotspots). This visualises directly why uniform
covers more: training commits probability mass to a few revisited cells. The shared log scale is
essential here---on the earlier per-panel linear scale these $1000\times$ dwell spikes saturated each
panel and hid the surrounding path structure.

\begin{figure}[h]
\centering
\includegraphics[width=0.6\textwidth]{../LPM_exploration/Miniworld/experiments/figures/fig_tv_fixation.png}
\caption{Share of steps spent on the noisy-TV action (action~4) under \texttt{action\_noise},
mean$\pm$std over 64 seeds. LPM $\approx3.0\%$; MSE $\approx83\%$; uniform $=1/5$ (dotted line).}
\label{fig:tv}
\end{figure}

\begin{figure}[h]
\centering
\includegraphics[width=\textwidth]{../LPM_exploration/Miniworld/experiments/figures/fig_coverage_curves.png}
\caption{Reachable-cell coverage vs.\ frames, per variant (mean$\pm$std over 64 seeds). The
uniform-random control is highest in every panel; all trained policies fall below it.}
\label{fig:curves}
\end{figure}

\begin{figure}[h]
\centering
\includegraphics[width=\textwidth]{../LPM_exploration/Miniworld/experiments/figures/fig_heatmap_evolution_action_noise_density.png}
\caption{Occupancy-density evolution (\texttt{action\_noise}), \emph{averaged over all 64 seeds}.
Rows: policy; columns: five equal training-time windows. Brightness is the mean per-cell visit count
on a \emph{shared log colour scale} (colorbar at right, $\sim\!10^{-1}$ to $\sim\!10^{3}$ visits);
the red dashed line is the noise wall ($z=8$), white rectangles are the reachable rooms. Uniform
spreads broadly at moderate counts; MSE/\textsc{none} show bright concentrated dwell cells.}
\label{fig:heat}
\end{figure}

\clearpage
\section{Infrastructure and reproducibility}
The upstream maze experiment exists only as three Jupyter notebooks with an embedded A2C and no CLI.
We extracted the A2C engine, the LPM prediction/uncertainty nets, and the pixel decoder
\emph{verbatim} into an importable, CLI-driven package
(\texttt{LPM\_exploration/Miniworld/experiments/}), aligned its hyperparameters to the paper's
Appendix~C.2, \textbf{corrected the LPM intrinsic reward to the paper's log-space Eq.~(3)} (with the
raw notebook form retained behind \texttt{reward\_space="raw"}), added canonical ICM and RND on the
shared CNN encoder, a \texttt{--random-policy} uniform control, and a per-step position log. The
correction and the accompanying error-model learning-rate fix are covered by unit tests
(\texttt{tests/test\_models.py}). The environment geometry is reused from our own
\texttt{miniworld\_play/envs.py} (a faithful port of the upstream maze). Each run writes a per-update
CSV (\texttt{coverage\_frac, beyond\_wall\_frac, time\_at\_wall\_frac,} plus model/policy losses) and
a compact \texttt{.npz} of $(\text{step}, x, z, \text{action}, \text{sticky})$. A grid runner
launches all runs in parallel (resumable); an analysis script aggregates the CSVs and position logs
into the coverage table, the coverage curves, the TV-fixation figure, and the
\emph{seed-averaged} heatmap-evolution figures. All deviations are logged in
\texttt{LPM\_exploration/UPSTREAM.md}.

\section{Limitations (stated up front)}
\begin{enumerate}[leftmargin=1.6em]
  \item \textbf{Single-environment variance.} We run the paper's budget ($5\times10^{4}$ steps) at
  $64$ seeds, but each run is still a \emph{single} environment, so coverage std stays at
  $0.08$--$0.32$; the closer rankings among trained methods are illustrative, not significance-tested.
  The uniform-beats-all and TV-fixation results, however, are large and robust. LPM specifically is
  \emph{bimodal} (explore vs.\ collapse), so its mean/median should be read together, not as a point
  estimate; a stability ablation (reward scaling / normalisation) is the natural follow-up.
  \item \textbf{Coverage is the wrong yardstick here.} In a small, extrinsic-reward-free maze,
  coverage rewards randomness, so it cannot demonstrate the value of directed exploration---which is
  why uniform wins. A faithful test of LPM's exploration benefit needs a sparse-reward,
  hard-exploration task (the paper's Atari / Montezuma settings), where no-intrinsic agents get zero
  reward. What the maze \emph{can} show---and does---is noise robustness, decoupled from coverage.
  \item \textbf{Figure-3 baselines are absent from the released maze code.} The maze notebooks contain
  only LPM and pixel-curiosity; AMA/EDT/EME/Ensemble/IDF/RND are not present, so a full reproduction of
  the paper's 7-method coverage ranking is not possible from this code. Our RND/ICM are our own ports.
  \item \textbf{Coverage is a 2D projection.} Occupancy is binned on the floor plane; agent heading
  and the first-person view (what actually drives the pixel predictors) are not visualised here.
\end{enumerate}

\end{document}
